% 1.3
% Ficheiro incluído por main.tex

\section[Promoção a Domain Controller (isec.local)]{Promoção a Domain Controller (\texttt{isec.local})}

\esquema{1-3-promocao-dc}{Domain Controller: AD, KDC Kerberos e DNS do domínio isec.local}

\passo{Passo 1 - Promover o servidor}

No Server Manager, clicar na \textbf{bandeira amarela} → \textbf{Promote this server to a domain controller}.

\passo{Passo 2 - Nova floresta}

\textbf{Add a new forest} → \emph{Root domain name:} \textbf{\texttt{isec.local}} → Next.

\passo{Passo 3 - Palavra-passe de restauro (DSRM)}

Definir a palavra-passe \textbf{DSRM} (\texttt{\textless{}PASSWORD\_LAB\textgreater{}}) e memorizá-la → Next.

\passo{Passo 4 - Aviso de DNS Delegation}

Ignorar o aviso de \textbf{DNS Delegation} e prosseguir com \textbf{Next}.

\passo{Passo 5 - Nome NetBIOS e instalação}

Verificar que o nome \textbf{NetBIOS} é \textbf{\texttt{ISEC}} → \textbf{Next} em todos os ecrãs seguintes → \textbf{Install}.

\passo{Passo 6 - Reinício e login no domínio}

O servidor reinicia sozinho para concluir o processo.

\emph{Verificação esperada:} no arranque seguinte, o login passa a ser feito como \textbf{\texttt{ISEC\textbackslash{}Administrator}} (Domain Administrator). Iniciar sessão e aguardar que o sistema arranque por completo.

\begin{simplesbox}
A promoção cria automaticamente a zona DNS \texttt{isec.local} e os registos SRV (\texttt{\_ldap.\_tcp.dc.\_msdcs.isec.local}, \texttt{\_kerberos.\_tcp.isec.local}) que o Kerberos e o SSSD usam para descobrir o DC.
\end{simplesbox}

\passo{Passo 7 - Forwarder DNS (acesso à internet a partir do cliente)}

\emph{Objetivo:} como o Ubuntu vai usar \textbf{só} o DNS do Windows, o DC tem de reencaminhar para fora os nomes que não conhece.

\begin{lstlisting}[style=psstyle]
Add-DnsServerForwarder -IPAddress 1.1.1.1
Get-DnsServerForwarder
\end{lstlisting}

\emph{Verificação esperada:} \texttt{Get-DnsServerForwarder} lista \texttt{1.1.1.1}.

\problemas{\hyperref[ec-a]{Erros Comuns - Parte A}}
